\chapter{The uniform reductions and their quantifiers}\label{ch:reductions}
\sourcebox{First helper prompt; rough-input and later proof notes}{These reductions are proved. The bounded-exponent or rough-input estimate to which they reduce the problem remains unproved for arbitrary targets.}

\section{Removing high input exponents}
Fix $E\ge1$. Split every preimage into coprime factors
\[
 k=D_EU_E,\qquad D_E=\prod_{v_p(k)>E}p^{v_p(k)}.
\]
The factor $U_E$ has all exponents at most $E$. Since $k\mid N$, the number
of possible high-exponent parts is at most
\[
 G_E(N)=\prod_{p\mid N}\left(1+\max\{\alpha_p-E,0\}\right).
\]
For each actual $D_E$, multiplicativity gives $h(U_E)=N/h(D_E)$. Therefore
\begin{equation}\label{red:fixedcap}
 \boxed{f(N)\le G_E(N)F_E(N).}
\end{equation}
Define
\[
 c_E=\sup_{a\in\mathbb Z,\ a\ge E+1}\frac{\log(a-E+1)}a.
\]
For fixed $E$, the same small/large-prime split $z=L/w^3$ as before gives
\begin{equation}\label{red:GE}
 \log G_E(N)\le(c_E+o_E(1))L/w.
\end{equation}
Indeed, the small primes contribute $O(L/w^2)$; for the large primes,
$\log(1+\max\{a-E,0\})\le c_Ea$ and $\sum_{p>z}\alpha_p\le L/\log z$.
For $E\ge2$,
\[
 0\le c_E\le\frac{\log(E+1)}{E+1}\longrightarrow0.
\]

\begin{theorem}\label{red:equivalence}
The original conjecture is equivalent to
$\log F_E(X)=o_E(\log X/\log\log X)$ for every fixed $E$.
\end{theorem}
\begin{proof}
The forward implication is immediate for sufficiently large targets; finitely
many smaller targets contribute a constant to the maximum. Conversely, given
$\varepsilon>0$, first fix $E$ so that $c_E<\varepsilon/3$. Next take $N$
large enough for the $o_E(1)$ in \eqref{red:GE} and the fixed-cap bound both
to contribute less than $\varepsilon/3$. Equation~\eqref{red:fixedcap} proves
the required inequality. There is no passage to $E=E(N)$ without additional
uniform error control.
\end{proof}
For $E=1$, $F_1(X)=1$ by squarefree injectivity and
$G_1(N)=\prod_{p\mid N}\alpha_p$. The framework recovers the starting bound.

\section{Also removing small input primes}
For fixed $E\ge1$ and $\delta>0$, define
\[
 R_{E,\delta}(X)=\max_{1\le M\le X}
 \#\{u:h(u)=M,\ v_p(u)\le E,\ P^-(u)>\delta\log X\}.
\]
If needed replace its logarithm by $\log\max\{1,R_{E,\delta}(X)\}$.
The cutoff is defined from the original $X$, not separately from each residual
$M$. After selecting $D_E$, fix the remaining exponents at primes at most
$\delta\log N$. This costs at most $(E+1)^{\pi(\delta\log N)}$ records.
The residual is counted by $R_{E,\delta}(N)$, so
\begin{equation}\label{red:rough}
 \boxed{f(N)\le G_E(N)(E+1)^{\pi(\delta\log N)}
 \max\{1,R_{E,\delta}(N)\}.}
\end{equation}
The prime number theorem gives
\[
 \pi(\delta\log N)\log(E+1)
 =\bigl(\delta\log(E+1)+o_{E,\delta}(1)\bigr)L/w.
\]
Thus the following family of estimates would suffice:
\begin{equation}\label{red:roughgap}
 \boxed{\log\max\{1,R_{E,\delta}(X)\}
 =o_{E,\delta}(\log X/\log\log X)
 \quad\text{for every fixed }E,\delta.}
\end{equation}
Choose $E$ first, then $\delta$, then the large target threshold. Neither
\eqref{red:roughgap} nor an unrestricted substitute is established in the record.

\section{A short interval, with its exact limitation}
If $h(u)=M\le X$ and $P^-(u)>y>1$, then
\[
 \log\frac{\sigma(u)}u
 \le\sum_{p\mid u}\log\frac p{p-1}
 \le\frac{\omega(u)}{y-1}
 \le\frac{\log X}{2(y-1)\log y}.
\]
Consequently,
\begin{equation}\label{red:interval}
 \sqrt M\exp\!\left(-\frac{\log X}{4(y-1)\log y}\right)
 \le u\le\sqrt M.
\end{equation}
For $y=\delta\log X$, this is relative width $O_\delta(1/\log\log X)$.
It does not count the divisors of $M$ satisfying the equation in that interval.
A separate arithmetic counting theorem is required. The nearby literature on
fibers of $\sigma(n)-n$ warns that fixed-relative-width concentration alone can
coexist with large fibers; it is not a theorem about this shrinking interval.

\section{Asymptotic bookkeeping checklist}
In any claimed completion, verify the maximum over all $M\le X$, including small
residuals; the dependence of constants on $E$ and $\delta$; the original cutoff
used in the roughness condition; and exact multiplicativity of the selected
unitary blocks. A proof at one fixed exponent cap gives partial progress, not
the all-cap implication. A numerical decrease of a fixed general constant does
not by itself imply its limit can be made zero.
